\noindent\textbf{Elección del procedimiento.} En estacionario el rotor sostiene $T=mg$. El disco ideal da $T=2\rho Av_i^2$. Despejar $v_i=\sqrt{T/(2\rho A)}$ y después calcular potencia inducida $P=Tv_i$.

\begin{enumerate}

\item El empuje compensa peso.

\[\fitmath{T=m\cdot g}\]

\[\fitmath{T=(5000)\cdot (9{,}81)\approx 49050\,\mathrm{N}}\]

\item El dato es radio del disco.

\[\fitmath{A=\pi\cdot {R}^{2}}\]

\[\fitmath{A=\pi\cdot {(8)}^{2}\approx 201{,}0619\,\mathrm{m^{2}}}\]

\item Tomar la raíz positiva de la relación de momento.

\[\fitmath{\mathrm{vi}=\sqrt{\frac{T}{2\cdot \rho\cdot A}}}\]

\[\fitmath{\mathrm{vi}=\sqrt{\frac{(49050)}{2\cdot (1{,}225)\cdot (201{,}0619298)}}\approx 9{,}978644\,\mathrm{m/s}}\]

\item La potencia ideal no incluye perfil de pala, transmisión ni rotor antipar.

\[\fitmath{P=T\cdot \mathrm{vi}}\]

\[\fitmath{P=(49050)\cdot (9{,}978644262)\approx 489452{,}5\,\mathrm{W}}\]

\end{enumerate}
